% --- INSTITUTIONAL TEMPLATE FED v7.4.9 ---
% Estándar de Reporteo APA 7 (Babel-Free)
\documentclass[12pt,spanish]{article}

\usepackage{graphicx}
\usepackage{geometry}
\usepackage{indentfirst} 
\usepackage{caption}
\usepackage{floatrow} 

% Forzar estilos de caption antes de cualquier carga
\floatsetup[figure]{capposition=top}
\floatsetup[table]{capposition=top}

\usepackage{booktabs}
\usepackage{float}
\usepackage{longtable}
\usepackage{threeparttable}
\usepackage{array}
\usepackage{calc}
\usepackage{amsmath, amssymb}
\usepackage{fancyhdr}
\usepackage{mathptmx} 
\usepackage[T1]{fontenc}
\usepackage{setspace} 
\usepackage{ragged2e} 
\usepackage{titlesec}
\usepackage{url}
\def\UrlBreaks{\do\/\do-\do\.\do\_\do\?\do\=\do\&\do\%}
\usepackage[hidelinks,breaklinks]{hyperref}

% --- Pandoc Compatibility Block (Iron Clad) ---
\providecommand{\tightlist}{%
  \setlength{\itemsep}{0pt}\setlength{\parskip}{0pt}}

\usepackage{xcolor}
\usepackage{fancyvrb}
\usepackage{framed}

% --- Code Highlighting (Pandoc Standard) ---
\AtBeginDocument{
  \definecolor{shadecolor}{RGB}{248,248,248}
}

% Compatibilidad Pandoc 3.x
\ifdefined\pandocbounded\else
  \newcommand{\pandocbounded}[1]{#1}
\fi

% --- Pandoc Citeproc (CSL) compatibility ---
\newlength{\cslhangindent}
\setlength{\cslhangindent}{1.27cm} 
\newlength{\csllabelwidth}
\setlength{\csllabelwidth}{3em}
\newcommand{\citeproctext}{} 
\newcommand{\citeprocitem}[2]{\item} 
\newenvironment{CSLReferences}[2] 
 {\begin{list}{}{
  \setlength{\leftmargin}{\cslhangindent}
  \setlength{\parsep}{0pt}
  \setlength{\itemsep}{6pt}
  \ifodd #1
   \setlength{\leftmargin}{\leftmargin + \csllabelwidth}
   \setlength{\itemindent}{-\csllabelwidth}
  \fi
  }
  \renewcommand{\bibitem}[2][]{\item} % Elimina los corchetes [ ]
 }
 {\end{list}}
\newcommand{\CSLBlock}[1]{#1\hfill\break}
\newcommand{\CSLLeftMargin}[1]{\parbox[t]{\csllabelwidth}{#1}}
\newcommand{\CSLRightInline}[1]{\parbox[t]{\linewidth - \csllabelwidth}{#1}\break}
\newcommand{\CSLIndent}[1]{\hspace{\cslhangindent}#1}

\makeatletter
\@ifundefined{paragraph}{\newcounter{paragraph}}{}
\@ifundefined{subparagraph}{\newcounter{subparagraph}}{}
\makeatother

% Solución al error 'No counter none defined' en longtable
\makeatletter
\newcounter{none}
\@ifpackageloaded{caption}{
  \DeclareCaptionType{none}
  \captionsetup[none]{name=Tabla}
}{}
\makeatother

\hypersetup{
  colorlinks=true,
  linkcolor=blue,
  filecolor=blue,
  citecolor=blue,
  urlcolor=blue
}

% --- Macros Institucionales ---
\newcommand{\fuente}[1]{
  \vspace{4pt}
  \begin{flushleft}
    \small
    \textit{Nota.} #1
  \end{flushleft}
}

% --- Localización Manual Determinista (Sin Babel) ---
\AtBeginDocument{
  \renewcommand{\tablename}{Tabla}
  \renewcommand{\figurename}{Figura}
  \renewcommand{\contentsname}{Índice}
  \renewcommand{\listtablename}{Índice de Tablas}
  \renewcommand{\listfigurename}{Índice de Figuras}
  \renewcommand{\abstractname}{Resumen}
  \renewcommand{\refname}{Referencias}
  
  \RaggedRight
  \setlength{\parindent}{1.27cm}
  \setlength{\parskip}{0pt}
}

\DeclareCaptionLabelSeparator{newline}{\\ \vspace{2pt}}
\captionsetup{
  position=top,
  justification=RaggedRight,
  singlelinecheck=false,
  labelfont=bf,
  textfont=it,
  labelsep=newline,
  font=small,
  skip=10pt
}
\captionsetup[figure]{name=Figura}
\captionsetup[table]{name=Tabla}

% Márgenes APA 7
\geometry{
  a4paper,
  margin=25.4mm,
  headheight=15mm
}

% --- Encabezados ---
\pagestyle{fancy}
\fancyhf{} 
\fancyhead[L]{\includegraphics[height=1.2cm]{/home/erick-fcs/Capital_Workstation/capital-workstation-libs/src/ecs_quantitative/reporting/assets/logos/logo_unl.png}}
\fancyhead[R]{\includegraphics[height=1.2cm]{/home/erick-fcs/Capital_Workstation/capital-workstation-libs/src/ecs_quantitative/reporting/assets/logos/logo_economia.jpg}}
\renewcommand{\headrulewidth}{0pt}
\fancyfoot[R]{\thepage}

\fancypagestyle{plain}{
  \fancyhf{}
  \fancyhead[L]{\includegraphics[height=1.2cm]{/home/erick-fcs/Capital_Workstation/capital-workstation-libs/src/ecs_quantitative/reporting/assets/logos/logo_unl.png}}
  \fancyhead[R]{\includegraphics[height=1.2cm]{/home/erick-fcs/Capital_Workstation/capital-workstation-libs/src/ecs_quantitative/reporting/assets/logos/logo_economia.jpg}}
  \fancyfoot[R]{\thepage}
  \renewcommand{\headrulewidth}{0pt}
}

% --- Títulos ---
\titleformat{\section}{\normalfont\normalsize\bfseries\centering}{\thesection}{1em}{}
\titleformat{\subsection}{\normalfont\normalsize\bfseries}{\thesubsection}{1em}{}
\titleformat{\subsubsection}{\normalfont\normalsize\bfseries\itshape}{\thesubsubsection}{1em}{}
\titlespacing*{\section}{0pt}{12pt}{6pt}
\titlespacing*{\subsection}{0pt}{12pt}{6pt}
\titlespacing*{\subsubsection}{0pt}{12pt}{6pt}

\newenvironment{referenciasAPA}{
  \setlength{\parindent}{-1.27cm}
  \setlength{\leftskip}{1.27cm}
  \setlength{\parskip}{6pt}
}{\par}

\makeatletter\@ifpackageloaded{caption}{}{\usepackage{caption}}\AtBeginDocument{%
\ifdefined\contentsname
  \renewcommand*\contentsname{Table of contents}
\else
  \newcommand\contentsname{Table of contents}
\fi
\ifdefined\listfigurename
  \renewcommand*\listfigurename{List of Figures}
\else
  \newcommand\listfigurename{List of Figures}
\fi
\ifdefined\listtablename
  \renewcommand*\listtablename{List of Tables}
\else
  \newcommand\listtablename{List of Tables}
\fi
\ifdefined\figurename
  \renewcommand*\figurename{Figure}
\else
  \newcommand\figurename{Figure}
\fi
\ifdefined\tablename
  \renewcommand*\tablename{Table}
\else
  \newcommand\tablename{Table}
\fi
}\@ifpackageloaded{float}{}{\usepackage{float}}
\floatstyle{ruled}
\@ifundefined{c@chapter}{\newfloat{codelisting}{h}{lop}}{\newfloat{codelisting}{h}{lop}[chapter]}
\floatname{codelisting}{Listing}\newcommand*\listoflistings{\listof{codelisting}{List of Listings}}\makeatother\makeatletter\makeatother\makeatletter\@ifpackageloaded{caption}{}{\usepackage{caption}}
\@ifpackageloaded{subcaption}{}{\usepackage{subcaption}}\makeatother

\begin{document}

\begin{center}
    {\bfseries \Large UNIVERSIDAD NACIONAL DE LOJA}\\[0.3cm]
    {\bfseries \large FACULTAD JURÍDICA, SOCIAL Y
ADMINISTRATIVA}\\[0.2cm]
    {\bfseries CARRERA DE ECONOMÍA}\\[1.5cm]
    
    {\bfseries \large POLÍTICA ECONÓMICA}\\[1cm]
    
    {\bfseries \Large Taller Aplicado: Unidad 1}\\[1.5cm]
\end{center}

\noindent {\bfseries Estudiante:} Erick Fabricio Condoy Seraquive \\
\noindent {\bfseries Docente:} Econ. Maria Gabriela Moreno Hurtado \\
\noindent {\bfseries Fecha:} 20 de abril de 2026 \\

\vspace{1cm}

\doublespacing
\setlength{\parindent}{1.27cm}

\section{Introducción al Caso
Andesur}\label{introducciuxf3n-al-caso-andesur}

El país enfrenta un escenario de crecimiento débil, malestar social y
presiones sectoriales divergentes. A continuación, se presenta el
análisis técnico bajo el marco de Cuadrado Roura (2014).

\section{Parte I: Conceptos
Fundamentales}\label{parte-i-conceptos-fundamentales}

\subsection{1. Economía vs.~Política
Económica}\label{economuxeda-vs.-poluxedtica-econuxf3mica}

\textbf{a) Economía:} Es el diagnóstico del caso (ej. el crecimiento ha
sido bajo en los últimos 3 años, el desempleo juvenil aumentó). Se
limita a describir y explicar los fenómenos. \textbf{b) Política
Económica:} Es el conjunto de medidas anunciadas por el gobierno (ej.
congelamiento de precios, aumento de salario mínimo). Es la intervención
deliberada. \textbf{Necesidad de la Política:} El caso no puede
entenderse solo desde la economía porque existe una \textbf{intervención
de poder}. Las demandas sociales y las presiones políticas exigen que el
Estado modifique las variables para alcanzar fines sociales que el
mercado por sí solo no está proveyendo.

\subsection{2. Praxis vs.~Disciplina}\label{praxis-vs.-disciplina}

\begin{itemize}
\tightlist
\item
  \textbf{Praxis (Práctica):} El anuncio del plan económico por parte
  del gobierno y la ejecución de las transferencias monetarias.
\item
  \textbf{Disciplina (Teoría):} El análisis de los economistas sobre la
  incoherencia del plan y la evaluación de la viabilidad fiscal por
  parte del Ministerio de Finanzas. \textbf{Riesgo de Confusión:}
  Confundir praxis con disciplina lleva a decisiones voluntaristas
  (creer que solo con desear un fin se tienen los instrumentos técnicos)
  o a análisis académicos alejados de la factibilidad política.
\end{itemize}

\section{Parte II: Análisis Positivo y
Normativo}\label{parte-ii-anuxe1lisis-positivo-y-normativo}

\subsection{3. Clasificación de
Enunciados}\label{clasificaciuxf3n-de-enunciados}

\begin{enumerate}
\def\labelenumi{\alph{enumi})}
\tightlist
\item
  \textbf{Positivo:} ``El desempleo juvenil aumentó\ldots{}'' (Hecho
  verificable).
\item
  \textbf{Normativo:} ``El gobierno debería priorizar\ldots{}'' (Juicio
  de valor/prioridad).
\item
  \textbf{Positivo:} ``El crédito subsidiado puede aumentar\ldots{}''
  (Relación causa-efecto).
\item
  \textbf{Normativo:} ``Es injusto que las familias\ldots{}'' (Carga
  moral/ética).
\item
  \textbf{Positivo:} ``La reducción de ciertos impuestos podría
  estimular\ldots{}'' (Hipótesis técnica).
\item
  \textbf{Normativo:} ``La prioridad debe ser el bienestar\ldots{}''
  (Criterio de valor político).
\end{enumerate}

\subsection{4. Relación en la
Práctica}\label{relaciuxf3n-en-la-pruxe1ctica}

Aunque se distinguen conceptualmente, se relacionan porque el análisis
\textbf{positivo} (el diagnóstico del bajo crecimiento) es el insumo
necesario para que el decisor aplique su criterio \textbf{normativo} (la
decisión de qué objetivo priorizar).

\section{Parte III: Jerarquía de la Política
Económica}\label{parte-iii-jerarquuxeda-de-la-poluxedtica-econuxf3mica}

\subsection{5. Tabla de Fines, Objetivos e
Instrumentos}\label{tabla-de-fines-objetivos-e-instrumentos}

{\def\LTcaptype{none} % do not increment counter
\begin{longtable}[]{@{}
  >{\raggedright\arraybackslash}p{(\linewidth - 4\tabcolsep) * \real{0.3333}}
  >{\raggedright\arraybackslash}p{(\linewidth - 4\tabcolsep) * \real{0.3333}}
  >{\raggedright\arraybackslash}p{(\linewidth - 4\tabcolsep) * \real{0.3333}}@{}}
\toprule\noalign{}
\begin{minipage}[b]{\linewidth}\raggedright
Fines (Metas Genéricas)
\end{minipage} & \begin{minipage}[b]{\linewidth}\raggedright
Objetivos (Metas Cuantificables)
\end{minipage} & \begin{minipage}[b]{\linewidth}\raggedright
Instrumentos (Medios)
\end{minipage} \\
\midrule\noalign{}
\endhead
\bottomrule\noalign{}
\endlastfoot
Bienestar de las familias & Reducción del desempleo juvenil & Crédito
subsidiado \\
Mayor equidad social & Estabilidad de precios & Reducción temporal de
impuestos \\
& & Inversión pública en obras \\
& & Transferencias monetarias \\
& & Revisión de aranceles \\
& & Aumento del salario mínimo \\
& & Formalización de pequeños negocios \\
\end{longtable}
}

\textbf{Nota sobre confusión:} No debe confundirse un instrumento con un
fin porque el instrumento es el ``cómo'' y el fin es el ``para qué''
último. Usar un instrumento (ej. aranceles) como si fuera un fin puede
llevar a políticas proteccionistas ineficientes.

\section{Parte IV: Coherencia y
Clasificación}\label{parte-iv-coherencia-y-clasificaciuxf3n}

\subsection{6. Análisis de Coherencia de
Instrumentos}\label{anuxe1lisis-de-coherencia-de-instrumentos}

\begin{enumerate}
\def\labelenumi{\arabic{enumi}.}
\tightlist
\item
  \textbf{Aumento Salario Mínimo:} Persigue el bienestar; podría
  aumentar el consumo; pero genera tensión al aumentar costos para
  microempresas (PYMES), pudiendo causar más desempleo.
\item
  \textbf{Congelamiento de Precios:} Persigue estabilidad; alivia el
  gasto del hogar; pero genera costos de escasez y desincentiva la
  oferta.
\item
  \textbf{Inversión Pública en Obras:} Persigue empleo; dinamiza la
  economía local; pero choca con la limitación del espacio fiscal.
\end{enumerate}

\subsection{7. Taxonomía de
Clasificación}\label{taxonomuxeda-de-clasificaciuxf3n}

{\def\LTcaptype{none} % do not increment counter
\begin{longtable}[]{@{}
  >{\raggedright\arraybackslash}p{(\linewidth - 8\tabcolsep) * \real{0.2000}}
  >{\raggedright\arraybackslash}p{(\linewidth - 8\tabcolsep) * \real{0.2000}}
  >{\raggedright\arraybackslash}p{(\linewidth - 8\tabcolsep) * \real{0.2000}}
  >{\raggedright\arraybackslash}p{(\linewidth - 8\tabcolsep) * \real{0.2000}}
  >{\raggedright\arraybackslash}p{(\linewidth - 8\tabcolsep) * \real{0.2000}}@{}}
\toprule\noalign{}
\begin{minipage}[b]{\linewidth}\raggedright
Medida
\end{minipage} & \begin{minipage}[b]{\linewidth}\raggedright
Orientación
\end{minipage} & \begin{minipage}[b]{\linewidth}\raggedright
Carácter Inst.
\end{minipage} & \begin{minipage}[b]{\linewidth}\raggedright
Nivel Actuación
\end{minipage} & \begin{minipage}[b]{\linewidth}\raggedright
Temporal
\end{minipage} \\
\midrule\noalign{}
\endhead
\bottomrule\noalign{}
\endlastfoot
IVA/Impuestos & Proceso & Cuantitativa & Micro/Macro & Corto Plazo \\
Salario Mínimo & Ordenación & Cualitativa/Inst. & Macro/Social & Mediano
Plazo \\
Transferencias & Proceso & Cuantitativa & Social/Micro & Corto Plazo \\
\end{longtable}
}

\subsection{8. Multi-categoría}\label{multi-categoruxeda}

Una medida como la \textbf{Inversión Pública} puede ser
\textbf{Cuantitativa} (monto invertido) y de \textbf{Proceso} (atender
desempleo actual), pero también \textbf{Cualitativa} si implica un
cambio en la ley de contrataciones. No es error, sino reflejo de la
complejidad económica.

\section{Parte V: Tensiones y
Recomendaciones}\label{parte-v-tensiones-y-recomendaciones}

\subsection{9. Tensiones Técnicas}\label{tensiones-tuxe9cnicas}

\begin{enumerate}
\def\labelenumi{\arabic{enumi}.}
\tightlist
\item
  \textbf{Salario Mínimo vs.~Formalización:} El aumento del salario
  podría expulsar a pequeños negocios hacia la informalidad al no poder
  costearlo.
\item
  \textbf{Inversión Pública vs.~Espacio Fiscal:} Existe una
  contradicción directa entre la demanda de obra pública y la
  advertencia de Finanzas sobre la falta de recursos.
\end{enumerate}

\subsection{10. Recomendaciones de
Implementación}\label{recomendaciones-de-implementaciuxf3n}

\textbf{Inmediatas (4):} 1. \textbf{Formalización de negocios:} Bajo
costo fiscal y mejora la base tributaria. 2. \textbf{Transferencias a
vulnerables:} Efecto alivio inmediato con focalización precisa. 3.
\textbf{Inversión pública (obras mano de obra):} Ataque directo al
desempleo. 4. \textbf{Mesa técnica de productividad:} Para alinear
finanzas y realidad sectorial.

\textbf{Postergadas (2):} 1. \textbf{Congelamiento de precios:} Genera
distorsiones de mercado difíciles de revertir. 2. \textbf{Crédito
subsidiado masivo:} Riesgo de impago y mayor presión al déficit si no
está bien fondeado.

\end{document}
